\documentclass{article}
\usepackage{graphicx} % Required for inserting images
\usepackage[english]{babel}
\usepackage{amsfonts}
\usepackage{amssymb}
\usepackage{amsmath}
\usepackage{a4wide}
\usepackage{xcolor}
\usepackage{enumitem}


\title{Stat 17 Week 3}
\author{Xiao Zhang}
\date{January 2026}

\begin{document}

\maketitle

\section*{Probability Topics — Part 1: Basics \& Counting}

\begin{itemize}
  \item \textbf{What is Probability}
  \begin{itemize}
    \item Probability describes the long-run proportion of times an outcome occurs
    \item Probabilities are always between 0 and 1
    \begin{itemize}
      \item $P=0$: impossible event
      \item $P=1$: certain event
    \end{itemize}
  \end{itemize}

  \item \textbf{Small Samples vs Large Samples}
  \begin{itemize}
    \item Results from small samples can be misleading
    \item Short-term randomness does not reflect long-term behavior
  \end{itemize}

  \item \textbf{Law of Large Numbers}
  \begin{itemize}
    \item As the number of trials increases, observed proportions approach theoretical probabilities
    \item Early imbalances tend to stabilize over time
  \end{itemize}

  \item \textbf{Independence and the Gambler's Fallacy}
  \begin{itemize}
    \item Trials are independent if previous outcomes do not affect future outcomes
    \item Past results do not change future probabilities
    \item The belief that an outcome is ``due'' is incorrect
  \end{itemize}

  \item \textbf{Probability Models and Sample Space}
  \begin{itemize}
    \item A probability model consists of:
    \begin{itemize}
      \item a list of all possible outcomes
      \item a probability assigned to each outcome
    \end{itemize}
    \item The sample space $S$ is the set of all possible outcomes
  \end{itemize}

  \item \textbf{Multi-step Experiments}
  \begin{itemize}
    \item Experiments with multiple stages can be represented using tree diagrams
    \item Example: two coin tosses $\rightarrow \{HH, HT, TH, TT\}$
  \end{itemize}

  \item \textbf{With Replacement vs Without Replacement}
  \begin{itemize}
    \item With replacement: the number of possible outcomes remains the same each trial
    \item Without replacement: the number of possible outcomes decreases after each trial
  \end{itemize}

  \item \textbf{Order of Selection}
  \begin{itemize}
    \item If order matters, different sequences count as different outcomes
    \item If order does not matter, only the selected elements matter
  \end{itemize}

  \item \textbf{Permutations (Order Matters)}
  \[
    nP_r = \frac{n!}{(n-r)!} = n(n-1) \cdots(n-r+1)
  \]

  \item \textbf{Combinations (Order Does Not Matter)}
  \[
    nC_r = \binom{n}{r} = \frac{n!}{r!(n-r)!} =  \frac{n(n-1) \cdots(n-r+1)}{r!}
  \]
\end{itemize}


\end{document}
